\noindent This experiment compared the $B-H$ responses of several candidate transformer-core materials when driven by the same AC current, with the aim of identifying the sample that simultaneously maximises relative permeability and minimises hysteresis loss. Each dataset was processed into a hysteresis loop, which then yielded maximum $\mu_{\mathrm{r}}$ and cycle loss estimates (Table~\ref{tab:summary-metrics}). Transformer iron clearly outperforms the other ferromagnets, achieving $\mu_{\mathrm{r,max}}\approx 3.3\times 10^{2}$ and only ${\sim}4.7\times 10^{5}\,\si{\watt\per\metre\cubed}$ of loss, while the CuNi alloy remains close to air ($\mu_{\mathrm{r}}\approx 1.3$) once the temperature is raised above \SI{40}{\celsius}, consistent with a transition through its Curie point.
$\\$
$\\$
The ideal transformer core therefore has a high local $\mu_{\mathrm{r}}$ to deliver large flux densities for modest $H$ fields and a narrow loop to keep the cycle loss low. Within the constraints of the apparatus, transformer iron best satisfies both criteria, whereas mild steel sacrifices efficiency through its wider hysteresis and CuNi lacks magnetic hardness or permeability for practical transformer operation.
$\\$
$\\$
Several systematic issues limit the absolute accuracy of the reported values: the integrator gain differed from its design by about \SI{40}{\percent}, the samples did not completely fill the measurement coil, and parasitic elements distort the waveform harmonics. These effects bias every dataset by a similar multiplicative factor, so the relative rankings reported here still reflect the intrinsic behaviour of the materials even if the magnitudes carry larger uncertainty.
$\\$
$\\$
Future work should therefore focus on refining the calibration loop (pure \SI{50}{\hertz} sinusoid with known amplitude), using samples that match the coil geometry, and extending the temperature/frequency matrix for CuNi to pin down its Curie-like transition with greater fidelity.
